% =========================
% L25.tex
% =========================

\chapter{L25 -- Sistemi non stazionari, stabilit\`a uniforme e Lyapunov tempo-variante}
\label{chap:L25}

\section{Chiusura del metodo indiretto e passaggio ai sistemi non stazionari}

Concludiamo anzitutto l'idea vista alla fine della lezione precedente sul \emph{metodo indiretto di Lyapunov}.

Quando il linearizzato di un sistema non lineare attorno a un equilibrio ha una parte stabile e una parte instabile, tramite un opportuno cambio di base si pu\`o scrivere, localmente, una dinamica del tipo
\[
\dot x = Ax + h(x),
\qquad
A=
\begin{pmatrix}
A_1 & 0\\
0 & A_2
\end{pmatrix},
\]
dove:
\begin{itemize}
    \item \(A_1\) raccoglie i modi con parte reale positiva,
    \item \(A_2\) raccoglie i modi con parte reale negativa,
    \item \(h(x)\) contiene i termini non lineari di ordine superiore.
\end{itemize}

Per studiare l'instabilit\`a si costruisce una funzione del tipo
\[
V(x)=x^\top P x,
\qquad
P=
\begin{pmatrix}
P_1 & 0\\
0 & -P_2
\end{pmatrix},
\]
con \(P_1>0\), \(P_2>0\).  
Il segno meno sul blocco \(P_2\) serve proprio a costruire una funzione che non sia positiva definita su tutto lo spazio, ma che sia utile per applicare un criterio di instabilit\`a alla Chetaev.

Lungo le traiettorie del sistema non lineare, per \(x\) piccolo, si ottiene una relazione del tipo
\[
L_fV(x)=x^\top x + 2x^\top P h(x),
\]
e il termine \(2x^\top P h(x)\) \`e di ordine superiore.  
Questo consente di concludere che, se nel linearizzato c'\`e almeno un modo esponenzialmente divergente, allora l'equilibrio del sistema non lineare \`e instabile.

\medskip

\noindent
A questo punto passiamo a una nuova classe di sistemi.

\section{Sistemi non stazionari}

Un sistema \`e detto \textbf{non stazionario} o \textbf{tempo-variante} se la dinamica dipende esplicitamente dal tempo:
\begin{equation}
\label{eq:L25-nonstat}
\dot x = f(x,t).
\end{equation}

Nel caso lineare si scrive
\begin{equation}
\label{eq:L25-ltv}
\dot x(t)=A(t)x(t).
\end{equation}

La differenza concettuale rispetto ai sistemi stazionari \`e fondamentale:
\begin{center}
\emph{per i sistemi tempo-varianti, conoscere gli autovalori istantanei di \(A(t)\) non basta pi\`u per concludere sulla stabilit\`a del sistema.}
\end{center}

\subsection{Primo controesempio: autovalori sempre negativi ma soluzione divergente}

Consideriamo il sistema
\[
\dot x(t)=A(t)x(t),
\qquad
A(t)=
\begin{pmatrix}
-1 & a(t)\\
0 & -1
\end{pmatrix},
\qquad
a(t)=e^{2t}.
\]
Gli autovalori di \(A(t)\) sono sempre
\[
\lambda_1(t)=\lambda_2(t)=-1.
\]
Se ragionassimo come nel caso stazionario, saremmo portati a dire che il sistema \`e stabile. Ma non \`e cos\`\i.

Infatti il sistema vale
\[
\begin{cases}
\dot x_1 = -x_1 + e^{2t}x_2,\\[4pt]
\dot x_2 = -x_2.
\end{cases}
\]
Dalla seconda equazione otteniamo immediatamente
\[
x_2(t)=e^{-t}x_2(0).
\]
Sostituendo nella prima:
\[
\dot x_1 + x_1 = e^{2t}e^{-t}x_2(0)=e^t x_2(0).
\]
Moltiplichiamo per il fattore integrante \(e^t\):
\[
\frac{d}{dt}\bigl(e^t x_1(t)\bigr)=e^{2t}x_2(0).
\]
Integrando:
\[
e^t x_1(t)=x_1(0)+\frac{x_2(0)}{2}(e^{2t}-1),
\]
quindi
\[
x_1(t)=e^{-t}x_1(0)+\frac{x_2(0)}{2}\bigl(e^t-e^{-t}\bigr).
\]
Se \(x_2(0)\neq 0\), allora \(x_1(t)\) cresce come \(e^t\), dunque diverge.

\medskip

\noindent
\textbf{Conclusione:}
nonostante gli autovalori istantanei siano sempre uguali a \(-1\), il sistema non \`e stabile.

\subsection{Messaggio da ricordare}

Per i sistemi non stazionari il criterio ``autovalori con parte reale negativa'' non vale pi\`u in generale.  
Conta l'evoluzione complessiva nel tempo, non lo spettro istantaneo della matrice \(A(t)\).

\IfFileExists{25_page-0001.jpg}{
\begin{figure}[h!]
\centering
\includegraphics[width=.82\textwidth]{25_page-0001.jpg}
\caption{Appunti originali: introduzione ai sistemi non stazionari e controesempi sullo spettro istantaneo.}
\end{figure}
}{}

\section{Stabilit\`a per sistemi dipendenti dal tempo}

Nel caso non stazionario la definizione di stabilit\`a deve tener conto del fatto che il comportamento pu\`o dipendere anche dall'istante iniziale \(t_0\).

Consideriamo il sistema \eqref{eq:L25-nonstat}, con equilibrio nell'origine.

\subsection{Definizione di stabilit\`a}

L'origine \(\bar x=0\) si dice \textbf{stabile} se
\[
\forall \varepsilon>0,\ \forall t_0\ge 0,\ \exists \delta=\delta(\varepsilon,t_0)>0
\]
tale che
\[
\|x(t_0)\|<\delta
\quad\Longrightarrow\quad
\|x(t)\|<\varepsilon
\qquad \forall t\ge t_0.
\]

La differenza rispetto al caso stazionario \`e che qui \(\delta\) pu\`o dipendere da \(t_0\).

\subsection{Definizione di stabilit\`a uniforme}

L'origine si dice \textbf{uniformemente stabile} se
\[
\forall \varepsilon>0,\ \exists \delta=\delta(\varepsilon)>0
\]
indipendente da \(t_0\), tale che
\[
\|x(t_0)\|<\delta
\quad\Longrightarrow\quad
\|x(t)\|<\varepsilon
\qquad \forall t\ge t_0.
\]

\medskip

\noindent
\textbf{Interpretazione.}
Nel caso uniforme, la soglia di sicurezza \(\delta\) dipende solo da quanto piccolo vogliamo l'errore finale (\(\varepsilon\)), ma non dipende dal tempo in cui facciamo partire il sistema.

\section{Funzioni di classe \(\mathcal K\) e \(\mathcal K_\infty\)}

Quando si lavora con Lyapunov nei sistemi non stazionari, si usano molto le cosiddette funzioni di classe \(\mathcal K\).

\subsection{Definizione}

Una funzione
\[
\alpha:[0,a)\to [0,+\infty)
\]
si dice di \textbf{classe \(\mathcal K\)} se:
\begin{itemize}
    \item \(\alpha\) \`e continua;
    \item \(\alpha\) \`e strettamente crescente;
    \item \(\alpha(0)=0\).
\end{itemize}

Se inoltre \(a=+\infty\) e
\[
\alpha(r)\to +\infty
\qquad \text{per } r\to+\infty,
\]
allora \(\alpha\) si dice di \textbf{classe \(\mathcal K_\infty\)}.

\medskip

\noindent
Queste funzioni servono per ``incastrare'' una funzione di Lyapunov tra due stime radiali dipendenti solo dalla norma dello stato.

\begin{figure}[h!]
\centering
\begin{tikzpicture}[scale=1]
\draw[->] (-0.2,0) -- (4.3,0) node[right] {$r$};
\draw[->] (0,-0.2) -- (0,3.4) node[above] {$\alpha(r)$};
\draw[thick,blue] (0,0) .. controls (0.7,0.05) and (1.2,0.8) .. (1.8,1.4)
.. controls (2.4,2.0) and (3.2,2.4) .. (4,3.0);
\node[blue] at (3.6,2.5) {\small classe \(\mathcal K\)};
\end{tikzpicture}
\caption{Andamento qualitativo di una funzione di classe \(\mathcal K\).}
\end{figure}

\section{Lemma fondamentale: una funzione positiva definita si stima con funzioni di classe \(\mathcal K\)}

Sia
\[
V:D\to\mathbb R,
\qquad
0\in D\subseteq \mathbb R^n,
\]
continua e positiva definita.  
Se \(B_r\subset D\), allora esistono due funzioni \(\alpha_1,\alpha_2\in\mathcal K\) tali che
\begin{equation}
\label{eq:L25-K-bounds}
\alpha_1(\|x\|)\le V(x)\le \alpha_2(\|x\|)
\qquad \forall x\in B_r.
\end{equation}

Se inoltre \(D=\mathbb R^n\) e \(V\) \`e radialmente illimitata, allora si possono scegliere
\[
\alpha_1,\alpha_2\in\mathcal K_\infty.
\]

\medskip

\noindent
\textbf{Interpretazione.}
La disuguaglianza \eqref{eq:L25-K-bounds} dice che una funzione di Lyapunov positiva definita si comporta come una funzione crescente della distanza dall'origine.

\section{Caratterizzazione della stabilit\`a uniforme}

Un lemma molto utile afferma che l'origine \`e uniformemente stabile se e solo se esistono:
\begin{itemize}
    \item una funzione \(\alpha\in\mathcal K\),
    \item una costante \(c>0\),
\end{itemize}
tali che
\[
\|x(t)\|\le \alpha\bigl(\|x(t_0)\|\bigr)
\qquad \forall t\ge t_0,
\qquad \forall \|x(t_0)\|<c.
\]

\medskip

\noindent
Questa formulazione \`e comoda perch\'e permette di tradurre la stabilit\`a uniforme in una stima quantitativa esplicita.

\IfFileExists{25_page-0002.jpg}{
\begin{figure}[h!]
\centering
\includegraphics[width=.82\textwidth]{25_page-0002.jpg}
\caption{Appunti originali: definizioni di stabilit\`a uniforme e funzioni di classe \(\mathcal K\).}
\end{figure}
}{}

\section{Teorema di Lyapunov per sistemi non stazionari}

Passiamo ora al corrispondente criterio di Lyapunov per i sistemi tempo-varianti.

Consideriamo
\[
\dot x = f(x,t),
\qquad
f(0,t)=0 \ \ \forall t\ge 0.
\]
Sia
\[
V:[0,+\infty)\times D \to \mathbb R
\]
una funzione continuamente differenziabile.

Supponiamo che esistano due funzioni positive definite \(W_1,W_2\) tali che
\begin{equation}
\label{eq:L25-V-bounds}
W_1(x)\le V(t,x)\le W_2(x)
\qquad \forall t\ge 0,\ \forall x\in D,
\end{equation}
e che inoltre
\begin{equation}
\label{eq:L25-Vdot-timevarying}
\frac{\partial V}{\partial t}(t,x)
+
\frac{\partial V}{\partial x}(t,x)\,f(x,t)
\le 0.
\end{equation}

Allora l'origine \`e \textbf{uniformemente stabile}.

\medskip

\noindent
\textbf{Commento.}
La differenza rispetto al caso autonomo \`e che adesso la derivata totale lungo le traiettorie diventa
\[
\frac{d}{dt}V(t,x(t))
=
\frac{\partial V}{\partial t}(t,x(t))
+
\frac{\partial V}{\partial x}(t,x(t))\,\dot x(t),
\]
perch\'e \(V\) dipende esplicitamente anche dal tempo.

Se invece vale la stima pi\`u forte
\[
\frac{\partial V}{\partial t}(t,x)
+
\frac{\partial V}{\partial x}(t,x)\,f(x,t)
\le -W_3(x),
\]
con \(W_3\) positiva definita, allora si ottiene una forma di stabilit\`a pi\`u forte; nel caso lineare, addirittura stabilit\`a uniforme esponenziale.

\IfFileExists{25_page-0003.jpg}{
\begin{figure}[h!]
\centering
\includegraphics[width=.82\textwidth]{25_page-0003.jpg}
\caption{Appunti originali: teorema di Lyapunov per sistemi non stazionari e interpretazione geometrica tramite insiemi di livello.}
\end{figure}
}{}

\section{Esempio: funzione di Lyapunov che dipende dal tempo}

Consideriamo il sistema
\begin{equation}
\label{eq:L25-esempio-g}
\begin{cases}
\dot x_1 = -x_1 - g(t)x_2,\\[4pt]
\dot x_2 = x_1 - x_2,
\end{cases}
\end{equation}
dove \(g(t)\) \`e continua, derivabile e soddisfa
\[
0\le g(t)\le k,
\qquad
\dot g(t)\le g(t)
\qquad \forall t\ge 0.
\]

Vogliamo mostrare che l'origine \`e uniformemente stabile.

\subsection{Scelta della candidata}

Una candidata naturale \`e
\[
V(t,x)=x_1^2 + (1+g(t))x_2^2.
\]
Poich\'e \(0\le g(t)\le k\), otteniamo le stime
\[
x_1^2+x_2^2 \le V(t,x)\le x_1^2+(1+k)x_2^2.
\]
Dunque possiamo porre
\[
W_1(x)=x_1^2+x_2^2,
\qquad
W_2(x)=x_1^2+(1+k)x_2^2.
\]

\subsection{Derivata lungo le traiettorie}

Calcoliamo
\[
\frac{d}{dt}V(t,x(t))
=
\frac{\partial V}{\partial t}
+
\frac{\partial V}{\partial x}\dot x.
\]
Si ha
\[
\frac{\partial V}{\partial t}=\dot g(t)x_2^2,
\qquad
\frac{\partial V}{\partial x}=(2x_1,\ 2(1+g(t))x_2).
\]
Quindi
\begin{align*}
\dot V
&=
\dot g(t)x_2^2
+2x_1(-x_1-gx_2)
+2(1+g)x_2(x_1-x_2)\\
&=
\dot g\,x_2^2 -2x_1^2 -2g x_1x_2 +2x_1x_2 +2g x_1x_2 -2x_2^2 -2g x_2^2\\
&=
-2x_1^2 +2x_1x_2 -(2+2g-\dot g)x_2^2.
\end{align*}
Poich\'e \(\dot g \le g\), segue
\[
2+2g-\dot g \ge 2+g \ge 1,
\]
e dunque
\[
\dot V \le -2x_1^2 + 2x_1x_2 - x_2^2.
\]
Ma
\[
-2x_1^2 + 2x_1x_2 - x_2^2
=
-(x_1^2 + (x_1-x_2)^2)\le 0.
\]
Quindi \(\dot V \le 0\) e, per il teorema precedente, l'origine \`e uniformemente stabile.

\medskip

\noindent
\textbf{Osservazione.}
Qui non basta una funzione \(V(x)\) dipendente solo dallo stato: la scelta giusta \`e una funzione \(V(t,x)\) che incorpora esplicitamente il coefficiente tempo-variante \(g(t)\).

\IfFileExists{25_page-0004.jpg}{
\begin{figure}[h!]
\centering
\includegraphics[width=.72\textwidth]{25_page-0004.jpg}
\caption{Calcolo svolto a mano per l'esempio con \(g(t)\).}
\end{figure}
}{}

\section{Caso lineare non stazionario e equazione differenziale di Lyapunov}

Consideriamo ora il caso lineare
\[
\dot x = A(t)x,
\]
con \(A(t)\) continua.

Supponiamo che esista una matrice simmetrica \(P(t)\), continua e derivabile, tale che
\begin{equation}
\label{eq:L25-P-bounds}
0<c_1 I \le P(t)\le c_2 I
\qquad \forall t\ge t_0,
\end{equation}
e che soddisfi l'equazione differenziale di Lyapunov
\begin{equation}
\label{eq:L25-lyap-diff}
-\dot P(t)=P(t)A(t)+A^\top(t)P(t)+Q(t),
\end{equation}
dove \(Q(t)=Q^\top(t)\) \`e continua e
\begin{equation}
\label{eq:L25-Q-lower}
Q(t)\ge c_3 I>0.
\end{equation}

Poniamo
\[
V(t,x)=x^\top P(t)x.
\]
Dalle stime \eqref{eq:L25-P-bounds} segue
\[
c_1\|x\|^2 \le V(t,x)\le c_2\|x\|^2.
\]
Inoltre
\begin{align*}
\dot V
&=
x^\top \dot P x + x^\top P \dot x + \dot x^\top P x\\
&=
x^\top\bigl(\dot P + PA + A^\top P\bigr)x\\
&=
-x^\top Q(t)x\\
&\le -c_3\|x\|^2.
\end{align*}
Usando \(V\le c_2\|x\|^2\), otteniamo
\[
\dot V \le -\frac{c_3}{c_2}V.
\]
Per confronto scalare,
\[
V(t)\le e^{-\frac{c_3}{c_2}(t-t_0)}V(t_0),
\]
e quindi
\[
\|x(t)\|
\le
\sqrt{\frac{c_2}{c_1}}
\,e^{-\frac{c_3}{2c_2}(t-t_0)}
\|x(t_0)\|.
\]

\medskip

\noindent
\textbf{Conclusione.}
L'origine \`e \textbf{uniformemente esponenzialmente stabile}.

\section{Richiami: LaSalle--Yoshizawa e lemma di Barb\u{a}lat}

Questi risultati vengono spesso usati per rafforzare i criteri di Lyapunov nei sistemi tempo-varianti.

\subsection{Principio di LaSalle--Yoshizawa}

Se esiste una funzione \(V(t,x)\) tale che
\[
\alpha_1(\|x\|)\le V(t,x)\le \alpha_2(\|x\|),
\]
e
\[
\frac{\partial V}{\partial t} + \frac{\partial V}{\partial x}f(x,t)\le -W(x)\le 0,
\]
con \(W\) semidefinita positiva, allora le traiettorie restano limitate e
\[
W(x(t))\to 0 \qquad \text{per } t\to+\infty.
\]
Se si riesce poi a capire quali traiettorie possono restare nell'insieme
\[
S=\{x:W(x)=0\},
\]
si pu\`o concludere verso quale insieme tende la soluzione.

\subsection{Lemma di Barb\u{a}lat}

Una forma pratica molto usata \`e la seguente: se una funzione \(r(t)\) \`e uniformemente continua e il suo integrale converge, allora
\[
r(t)\to 0
\qquad \text{per } t\to+\infty.
\]

In applicazioni Lyapunov, spesso si prende
\[
r(t)=\dot V(t,x(t)),
\]
si mostra che \(\dot V\le 0\), che \(V\) \`e inferiormente limitata e che \(\dot V\) \`e uniformemente continua; da qui si conclude
\[
\dot V(t)\to 0.
\]

\IfFileExists{25_page-0005.jpg}{
\begin{figure}[h!]
\centering
\includegraphics[width=.82\textwidth]{25_page-0005.jpg}
\caption{Appunti originali: richiamo a LaSalle--Yoshizawa e al lemma di Barb\u{a}lat.}
\end{figure}
}{}

\section{Esercizio d'esame: 13/04/2021}
\label{sec:L25-esercizio-130421}

Consideriamo il sistema
\begin{equation}
\label{eq:L25-ex-130421}
\begin{cases}
\dot x_1 = x_1^2 - x_1x_2,\\[4pt]
\dot x_2 = -x_2^2 + a x_1x_2,
\end{cases}
\qquad a\in\mathbb R.
\end{equation}

\subsection{Punto 1: equilibri e metodo indiretto}

Gli equilibri si ottengono imponendo
\[
\begin{cases}
0 = x_1^2 - x_1x_2 = x_1(x_1-x_2),\\[4pt]
0 = -x_2^2 + a x_1x_2 = x_2(ax_1-x_2).
\end{cases}
\]

Dalla prima equazione:
\[
x_1=0
\qquad \text{oppure} \qquad
x_1=x_2.
\]

\paragraph{Caso A: \(x_1=0\).}
La seconda equazione diventa
\[
0=-x_2^2,
\]
quindi necessariamente \(x_2=0\).  
Otteniamo l'equilibrio
\[
(0,0).
\]

\paragraph{Caso B: \(x_1=x_2\).}
La seconda equazione diventa
\[
0=(a-1)x_2^2.
\]
Quindi:
\begin{itemize}
    \item se \(a\neq 1\), allora \(x_2=0\) e si torna a \((0,0)\);
    \item se \(a=1\), allora ogni punto della retta
    \[
    x_1=x_2
    \]
    \`e equilibrio.
\end{itemize}

\medskip

\noindent
\textbf{Conclusione sugli equilibri.}
\begin{itemize}
    \item Se \(a\neq 1\), l'unico equilibrio \`e \((0,0)\).
    \item Se \(a=1\), c'\`e una retta di punti di equilibrio:
    \[
    \{(s,s):s\in\mathbb R\}.
    \]
\end{itemize}

\subsection{Linearizzazione}

La matrice jacobiana \`e
\[
A(x)=
\begin{pmatrix}
2x_1-x_2 & -x_1\\
a x_2 & -2x_2 + a x_1
\end{pmatrix}.
\]

\paragraph{All'origine.}
\[
A(0,0)=
\begin{pmatrix}
0 & 0\\
0 & 0
\end{pmatrix}.
\]
Tutti gli autovalori sono nulli, quindi il metodo indiretto non conclude nulla.

\paragraph{Sulla retta di equilibrio, quando \(a=1\).}
Per \(x_1=x_2=s\),
\[
A(s,s)=
\begin{pmatrix}
s & -s\\
s & -s
\end{pmatrix}.
\]
Si ha
\[
\operatorname{tr}A=0,
\qquad
\det A=0,
\]
quindi gli autovalori sono ancora entrambi nulli. Anche qui il metodo indiretto \`e inconcludente.

\subsection{Punto 2: si pu\`o stabilizzare con un solo ingresso?}

Supponiamo di aggiungere un solo ingresso in una sola equazione. Ad esempio:
\[
\begin{cases}
\dot x_1 = x_1^2 - x_1x_2 + u,\\[4pt]
\dot x_2 = -x_2^2 + a x_1x_2,
\end{cases}
\]
oppure, analogamente, nella seconda.

Se linearizziamo all'origine, otteniamo sempre
\[
A(0,0)=0.
\]
Nel primo caso
\[
B=
\begin{pmatrix}
1\\
0
\end{pmatrix},
\]
nel secondo
\[
B=
\begin{pmatrix}
0\\
1
\end{pmatrix}.
\]
In entrambi i casi la matrice di raggiungibilit\`a \`e
\[
\mathcal R=[B\ AB]=[B\ 0],
\]
e quindi ha rango \(1<2\).  
Dunque la coppia \((A,B)\) non \`e completamente raggiungibile e il sistema non \`e completamente stabilizzabile con un solo ingresso locale all'origine.

\subsection{Punto 3: studio dell'instabilit\`a con Chetaev nel caso \(a=1\)}

Nel caso \(a=1\), proviamo la candidata
\[
V(x)=x_1^2-x_2^2.
\]
Questa funzione \emph{non} \`e definita positiva, quindi non serve per dimostrare stabilit\`a.  
Pu\`o per\`o essere usata in senso ``inverso'' per dimostrare instabilit\`a.

Calcoliamo la derivata lungo le traiettorie:
\begin{align*}
L_fV
&= 2x_1(x_1^2-x_1x_2) -2x_2(-x_2^2+x_1x_2)\\
&=2x_1^3-2x_1^2x_2+2x_2^3-2x_1x_2^2\\
&=2(x_1-x_2)^2(x_1+x_2).
\end{align*}

Consideriamo la regione
\[
U=\left\{x:\ x_1^2+x_2^2<\rho^2,\ \ x_1^2-x_2^2>0,\ \ x_1+x_2>0\right\}.
\]
In \(U\):
\begin{itemize}
    \item \(V(x)>0\),
    \item \(L_fV(x)>0\),
    \item l'origine sta sul bordo di \(U\).
\end{itemize}
Per il teorema di Chetaev, l'origine \`e instabile.

\medskip

\noindent
\textbf{Osservazione importante.}
Questo \`e coerente con il fatto che per \(a=1\) l'origine non \`e un equilibrio isolato, ma appartiene a una retta di equilibri.

\IfFileExists{25_page-0006.jpg}{
\begin{figure}[h!]
\centering
\includegraphics[width=.76\textwidth]{25_page-0006.jpg}
\caption{Appunti originali sull'esercizio del 13/04/2021: equilibri, linearizzazione e commenti sulla stabilizzabilit\`a.}
\end{figure}
}{}

\IfFileExists{25_page-0007.jpg}{
\begin{figure}[h!]
\centering
\includegraphics[width=.70\textwidth]{25_page-0007.jpg}
\caption{Appunti originali: uso di una candidata non definita in segno per dimostrare instabilit\`a alla Chetaev.}
\end{figure}
}{}

\section{Esercizio d'esame: 07/04/2022}
\label{sec:L25-esercizio-070422}

Consideriamo il sistema discreto
\begin{equation}
\label{eq:L25-ex-070422}
\begin{cases}
x_1(k+1)= -\dfrac13 x_1(k)+\alpha x_1(k)^2 + 2u(k)x_2(k),\\[8pt]
x_2(k+1)= 2x_1(k)x_2(k)+\dfrac13 x_2(k).
\end{cases}
\end{equation}

\subsection{Punto 1: equilibri del sistema con \(u=0\)}

Poniamo \(u=0\). Gli equilibri \((\bar x_1,\bar x_2)\) soddisfano
\[
\begin{cases}
\bar x_1 = -\dfrac13 \bar x_1 + \alpha \bar x_1^2,\\[8pt]
\bar x_2 = 2\bar x_1\bar x_2 + \dfrac13 \bar x_2.
\end{cases}
\]

La prima equazione diventa
\[
\frac43 \bar x_1 - \alpha \bar x_1^2 = 0
\quad\Longleftrightarrow\quad
\bar x_1\left(\frac43 - \alpha \bar x_1\right)=0.
\]
Quindi
\[
\bar x_1=0
\qquad \text{oppure} \qquad
\bar x_1=\frac{4}{3\alpha}\quad (\alpha\neq 0).
\]

La seconda diventa
\[
\bar x_2 - 2\bar x_1\bar x_2 - \frac13 \bar x_2 = 0
\quad\Longleftrightarrow\quad
\bar x_2\left(\frac23 - 2\bar x_1\right)=0,
\]
cio\`e
\[
\bar x_2=0
\qquad \text{oppure} \qquad
\bar x_1=\frac13.
\]

\paragraph{Combinando i casi:}
\begin{itemize}
    \item \((0,0)\) \`e equilibrio per ogni \(\alpha\in\mathbb R\);
    \item se \(\alpha\neq 0\), esiste anche l'equilibrio
    \[
    \left(\frac{4}{3\alpha},\,0\right);
    \]
    \item se \(\alpha=4\), allora \(\bar x_1=\frac13\) e quindi tutti i punti
    \[
    \left(\frac13,\bar x_2\right),\qquad \bar x_2\in\mathbb R,
    \]
    sono equilibri.
\end{itemize}

\subsection{Stabilit\`a locale degli equilibri}

Con \(u=0\), il jacobiano rispetto allo stato \`e
\[
A(x)=
\begin{pmatrix}
-\dfrac13 + 2\alpha x_1 & 0\\[8pt]
2x_2 & 2x_1+\dfrac13
\end{pmatrix}.
\]

\paragraph{All'origine.}
\[
A(0,0)=
\begin{pmatrix}
-\dfrac13 & 0\\[6pt]
0 & \dfrac13
\end{pmatrix}.
\]
Attenzione: qui bisogna essere coerenti con la forma del sistema. Se la seconda equazione \`e davvero
\[
x_2(k+1)=2x_1(k)x_2(k)+\frac13 x_2(k),
\]
allora gli autovalori all'origine sono
\[
\lambda_1=-\frac13,
\qquad
\lambda_2=\frac13.
\]
Poich\'e entrambi hanno modulo minore di \(1\), l'origine \`e localmente asintoticamente stabile.

\paragraph{All'equilibrio \(\left(\frac{4}{3\alpha},0\right)\), \(\alpha\neq 0\).}
Si ottiene
\[
A\!\left(\frac{4}{3\alpha},0\right)=
\begin{pmatrix}
\frac73 & 0\\[6pt]
0 & \frac{8+\alpha}{3\alpha}
\end{pmatrix}.
\]
Il primo autovalore vale sempre
\[
\lambda_1=\frac73,
\]
quindi ha modulo maggiore di \(1\). Ne segue che tale equilibrio \`e sempre instabile.

\paragraph{Sulla retta \(\left(\frac13,\bar x_2\right)\), quando \(\alpha=4\).}
In questo caso
\[
A\!\left(\frac13,\bar x_2\right)=
\begin{pmatrix}
\frac73 & 0\\[6pt]
2\bar x_2 & 1
\end{pmatrix}.
\]
Gli autovalori sono
\[
\lambda_1=\frac73,
\qquad
\lambda_2=1.
\]
Anche qui l'equilibrio non \`e asintoticamente stabile; in effetti \`e instabile perch\'e \(|\lambda_1|>1\).

\subsection{Punto 2: raggiungibilit\`a del linearizzato}

Il vettore d'ingresso \`e
\[
B(x)=
\begin{pmatrix}
\frac{\partial f_1}{\partial u}\\[4pt]
\frac{\partial f_2}{\partial u}
\end{pmatrix}
=
\begin{pmatrix}
2x_2\\
0
\end{pmatrix}.
\]

\paragraph{All'origine.}
\[
B(0,0)=
\begin{pmatrix}
0\\
0
\end{pmatrix},
\]
quindi il linearizzato non \`e controllabile.

\paragraph{All'equilibrio \(\left(\frac{4}{3\alpha},0\right)\).}
Anche qui
\[
B\!\left(\frac{4}{3\alpha},0\right)=
\begin{pmatrix}
0\\
0
\end{pmatrix},
\]
quindi ancora nessuna controllabilit\`a.

\paragraph{Sulla retta \(\left(\frac13,\bar x_2\right)\), con \(\alpha=4\).}
Se \(\bar x_2\neq 0\),
\[
B\!\left(\frac13,\bar x_2\right)=
\begin{pmatrix}
2\bar x_2\\
0
\end{pmatrix}.
\]
La matrice di raggiungibilit\`a \`e
\[
\mathcal R=
\begin{bmatrix}
B & AB
\end{bmatrix}
=
\begin{bmatrix}
2\bar x_2 & \dfrac{14}{3}\bar x_2\\[6pt]
0 & 4\bar x_2^2
\end{bmatrix}.
\]
Se \(\bar x_2\neq 0\), allora
\[
\operatorname{rank}\mathcal R = 2.
\]
Quindi il linearizzato \`e completamente raggiungibile e si pu\`o progettare un controllo locale stabilizzante.

\subsection{Forma di un controllo affine locale}

Se \((\bar x,\bar u)\) \`e il punto di equilibrio scelto, un controllo affine locale si scrive come
\[
u = K(x-\bar x)+\bar u.
\]
Esplicitamente, in dimensione \(2\),
\[
u = k_1(x_1-\bar x_1)+k_2(x_2-\bar x_2)+\bar u.
\]
Equivalentemente,
\[
u = k_1x_1+k_2x_2+k_3,
\]
dove \(k_3\) \`e scelto in modo da non spostare l'equilibrio.

\medskip

\noindent
Il progetto si fa come sempre sul sistema delle variazioni
\[
z(k+1)=Az(k)+Bv(k),
\qquad
v(k)=Kz(k),
\]
scegliendo \(K\) in modo che gli autovalori di \(A+BK\) cadano dove desiderato.

\IfFileExists{25_page-0008.jpg}{
\begin{figure}[h!]
\centering
\includegraphics[width=.76\textwidth]{25_page-0008.jpg}
\caption{Appunti originali sull'esercizio del 07/04/2022: equilibri e studio del linearizzato.}
\end{figure}
}{}

\IfFileExists{25_page-0009.jpg}{
\begin{figure}[h!]
\centering
\includegraphics[width=.76\textwidth]{25_page-0009.jpg}
\caption{Appunti originali: matrice di ingresso, raggiungibilit\`a e forma del controllo affine.}
\end{figure}
}{}

\section{Schema riassuntivo della lezione}

La lezione contiene quattro idee fondamentali.

\subsection*{1. Nei sistemi non stazionari la stabilit\`a dipende anche da \(t_0\)}

Per questo introduciamo:
\begin{itemize}
    \item stabilit\`a,
    \item stabilit\`a uniforme,
    \item funzioni di classe \(\mathcal K\).
\end{itemize}

\subsection*{2. Lyapunov si estende ai sistemi tempo-varianti}

La funzione candidata pu\`o dipendere da \(t\):
\[
V=V(t,x).
\]
La derivata corretta lungo le traiettorie \`e
\[
\frac{d}{dt}V(t,x(t))
=
\frac{\partial V}{\partial t}
+
\frac{\partial V}{\partial x}f(x,t).
\]

\subsection*{3. Nel caso lineare tempo-variante si usa una matrice \(P(t)\)}

Si ottiene l'equazione differenziale di Lyapunov
\[
-\dot P = PA + A^\top P + Q,
\]
che gioca lo stesso ruolo dell'equazione algebrica nel caso stazionario.

\subsection*{4. Gli esercizi d'esame combinano sempre le stesse idee}

In pratica, nei compiti bisogna saper fare bene queste quattro mosse:
\begin{enumerate}
    \item trovare gli equilibri;
    \item linearizzare correttamente;
    \item capire quando il metodo indiretto conclude e quando no;
    \item se serve, usare Chetaev, Lyapunov diretto o un controllo affine locale.
\end{enumerate}

\medskip

\noindent
\textbf{Frase da ricordare all'orale.}
\begin{center}
\emph{Per i sistemi tempo-varianti il criterio spettrale istantaneo non basta. La stabilit\`a si studia con strumenti pi\`u robusti, in particolare con funzioni di Lyapunov dipendenti dal tempo e con stime uniformi rispetto all'istante iniziale.}
\end{center}